% =======================================================
% 3. MODELLO DI SVILUPPO
% =======================================================
\section{Modello di Sviluppo}

\subsection{Scelta del modello}

Il gruppo adotta un approccio organizzativo iterativo-adattivo ispirato a Scrum,
basato su sprint settimanali con cadenza da martedì a martedì e rotazione
periodica dei ruoli.

La durata settimanale degli sprint consente di suddividere le attività in task
a granularità ridotta e maggiormente monitorabili, favorendo un controllo
continuo dello stato di avanzamento del progetto.

L'approccio adottato non viene inteso come modello incrementale in senso stretto:
durante l'avanzamento del progetto gli artefatti possono essere raffinati,
corretti, riorganizzati o parzialmente riscritti sulla base delle verifiche
interne, dei feedback ricevuti e delle criticità emerse. Tale scelta risulta più
adatta alla natura del progetto, poiché consente di recepire correzioni e di
migliorare il way of working\textsubscript{G}, la pianificazione
e la qualità degli artefatti prodotti sulla base delle evidenze raccolte sprint per sprint.

In fase RTB\textsubscript{G}, l'approccio supporta principalmente le attività di
analisi dei requisiti, analisi dello stato dell'arte, progettazione architetturale
preliminare e consolidamento della documentazione. A partire dallo Sprint 8, con
avvio previsto il 2026-05-26, viene esteso alle attività di setup del Proof of
Concept\textsubscript{G}. A partire dallo Sprint 9, tali attività proseguono con
l'implementazione delle funzionalità dimostrative previste per la Technology
Baseline\textsubscript{G}. A partire dal termine della Revisione RTB
(2026-06-23), lo stesso approccio iterativo-adattivo viene mantenuto per la fase
PB\textsubscript{G}, dove supporta l'implementazione del Minimum Viable
Product\textsubscript{G}, l'introduzione della pipeline di CI\textsubscript{G} e
della relativa suite di test automatizzati, e la redazione dei documenti tecnici
di dettaglio (Specifica Tecnica, Manuale Utente, Manuale Sviluppatore) previsti
dal regolamento di progetto per la Product Baseline.

\subsection{Motivazioni}

\begin{itemize}
    \item \textbf{Adattabilità:} il piano di lavoro può essere revisionato sprint
    per sprint in base ai risultati ottenuti, ai feedback della proponente e alle
    criticità emerse.

    \item \textbf{Visibilità continua:} le cerimonie\textsubscript{G} strutturate,
    gli strumenti di tracciamento adottati dal gruppo (GitHub Projects,
    issue\textsubscript{G}, branch\textsubscript{G} e pull request\textsubscript{G})
    e gli aggiornamenti asincroni permettono al gruppo di mantenere un quadro
    aggiornato dello stato di avanzamento.

    \item \textbf{Controllo periodico della qualità:} la Definition of
    Done\textsubscript{G} (Sez.~\ref{sec:dod}) e le attività di verifica riducono
    il rischio di accumulare errori, non conformità\textsubscript{G} o debito
    tecnico\textsubscript{G}.

    \item \textbf{Raffinamento iterativo degli artefatti:} i risultati
    dell'Analisi dei Requisiti, dell'Analisi dello Stato dell'Arte, del Proof of
    Concept\textsubscript{G} e delle retrospettive vengono utilizzati per
    correggere, consolidare e migliorare le scelte documentali, tecniche e
    organizzative.

    \item \textbf{Sviluppo di competenze trasversali:} la rotazione dei ruoli
    consente ai membri del gruppo di acquisire esperienza su aspetti gestionali,
    documentali, tecnici e di verifica.
\end{itemize}

\subsection{Cerimonie e cadenza settimanale}

Ogni sprint ha durata di una settimana (martedì~$\rightarrow$~martedì).
La riunione ordinaria del martedì si tiene alle \textbf{14:30} e segue un ordine
fisso coerente con il ciclo di controllo del Piano di Progetto:

\begin{enumerate}
    \item \textbf{Consuntivo di periodo e Sprint Review\textsubscript{G} dello
    sprint precedente:} verifica degli output prodotti rispetto agli obiettivi
    pianificati, controllo dei task completati o rimandati, rendicontazione delle
    attività svolte, delle ore effettive, dei costi sostenuti e degli scostamenti
    rispetto al preventivo.

    \item \textbf{Sprint Retrospective\textsubscript{G}:} analisi del grado di
    raggiungimento degli obiettivi previsti, individuazione delle cause degli
    scostamenti e valutazione delle criticità emerse nel processo di lavoro.

    \item \textbf{Misure correttive e aggiornamento dei rischi:} definizione delle
    azioni correttive da adottare, aggiornamento dell'analisi dei rischi e apertura
    di eventuali issue\textsubscript{G} di miglioramento.

    \item \textbf{Miglioramento della pianificazione futura:} revisione della
    pianificazione a breve termine e, se necessario, della pianificazione a lungo
    termine, sulla base dei risultati del consuntivo e della retrospettiva.

    \item \textbf{Aggiornamento del Preventivo a Finire:} aggiornamento della stima
    delle risorse residue e del costo stimato per completare il progetto.

    \item \textbf{Sprint Planning\textsubscript{G}:} definizione degli obiettivi
    dello sprint corrente, assegnazione dei ruoli e dei task, creazione delle
    issue e stima dei tempi previsti.
\end{enumerate}

Le decisioni, le azioni e le assegnazioni emerse durante la riunione vengono
registrate nel documento di riunione condiviso e successivamente formalizzate
nel verbale interno.

\begin{table}[H]
\centering
\renewcommand{\arraystretch}{1.3}
\begin{tabularx}{\textwidth}{|p{2.8cm}|p{2.2cm}|c|p{2cm}|X|}
    \hline
    \textbf{Momento} & \textbf{Giorno/Ora} & \textbf{Tipo} &
    \textbf{Durata} & \textbf{Descrizione} \\
    \hline
    Consuntivo di periodo e Sprint Review
        & Martedì, 14:30 & Sincrono & Variabile
        & Verifica degli output prodotti rispetto agli obiettivi pianificati,
          controllo dei task completati o rimandati, rendicontazione di ore
          effettive, costi sostenuti e scostamenti rispetto al preventivo. \\
    \hline
    Sprint Retrospective
        & Martedì, dopo il consuntivo & Sincrono & Variabile
        & Analisi del grado di raggiungimento degli obiettivi, individuazione
          delle cause degli scostamenti e valutazione delle criticità emerse. \\
    \hline
    Misure correttive e rischi
        & Martedì, dopo la retrospettiva & Sincrono & Variabile
        & Definizione delle azioni correttive, aggiornamento dell'analisi dei
          rischi e apertura di eventuali issue di miglioramento. \\
    \hline
    Aggiornamento pianificazione
        & Martedì & Sincrono & Variabile
        & Revisione della pianificazione a breve e lungo termine sulla base degli
          esiti del consuntivo, della retrospettiva e delle misure correttive. \\
    \hline
    Aggiornamento PaF
        & Martedì & Sincrono & Variabile
        & Aggiornamento del Preventivo a Finire, delle risorse residue e del costo
          stimato per completare il progetto. \\
    \hline
    Sprint Planning
        & Martedì, a chiusura del ciclo & Sincrono & Variabile
        & Definizione degli obiettivi dello sprint corrente, assegnazione di ruoli
          e task, creazione delle issue tramite Google Form\textsubscript{G} e
          stima dei tempi previsti. \\
    \hline
    Allineamento mid-sprint
        & Venerdì, 14:30 & Sincrono & $\sim$30 min
        & Incontro intermedio per risolvere dubbi, affrontare criticità operative,
          aggiornare il gruppo sullo stato delle attività e recepire eventuali
          osservazioni emerse dal Diario di Bordo\textsubscript{G}. \\
    \hline
    Daily Stand-up\textsubscript{G}
        & Durante lo sprint & Asincrono & $\sim$5 min
        & Aggiornamento asincrono utilizzato quando necessario per segnalare
          impedimenti, richieste di supporto, dubbi operativi o informazioni
          utili al coordinamento del gruppo. \\
    \hline
\end{tabularx}
\caption{Schema delle attività di controllo e pianificazione settimanale}
\end{table}

\subsection{Risorse Umane}

Il team è composto da 6 membri, con ruoli rotativi secondo quanto pianificato
nel presente Piano di Progetto. La rotazione dei ruoli e la distribuzione delle
ore sono monitorate sprint per sprint, con l'obiettivo di mantenere coerenza con
il preventivo iniziale e con gli impegni assunti in fase di candidatura.

Il monitoraggio operativo dell'avanzamento delle attività avviene tramite
GitHub Projects\textsubscript{G}, utilizzato per la gestione del backlog,
l'assegnazione delle issue e il controllo dello stato di avanzamento dello
sprint. Le informazioni di consuntivazione sono integrate mediante i fogli
condivisi utilizzati dal gruppo.
